
\begin{abstract}

Quantization can accelerate large language model (LLM) inference. Going beyond \texttt{INT8} quantization, the research community is actively exploring even lower precision, such as \texttt{INT4}. Nonetheless, state-of-the-art \texttt{INT4} quantization techniques only accelerate low-batch, edge LLM inference, failing to deliver performance gains in large-batch, cloud-based LLM serving. We uncover a critical issue: existing \texttt{INT4} quantization methods suffer from significant runtime overhead (20-90\%) when dequantizing either weights or partial sums on GPUs. To address this challenge, we introduce QoQ, a \texttt{W4A8KV4} quantization algorithm with 4-bit weight, 8-bit activation, and 4-bit KV cache. \algo stands for \textit{quattuor-octō-quattuor}, which represents 4-8-4 in Latin. QoQ is implemented by the \system inference library that achieves measured speedup. The key insight driving  \system is that the efficiency of LLM serving on GPUs is critically influenced by operations on low-throughput CUDA cores. Building upon this insight, in \algo algorithm, we introduce \textit{progressive quantization} that can allow low dequantization overhead in \texttt{W4A8} GEMM. Additionally, we develop \textit{SmoothAttention} to effectively mitigate the accuracy degradation incurred by 4-bit KV quantization. In the \system system, we perform \textit{compute-aware weight reordering} and take advantage of \textit{register-level parallelism} to reduce dequantization latency. We also transfer theoretical memory saving brought by KV4 attention into measured speedup using \system. 
As a result, \system improves the maximum achievable serving throughput of Llama-3-8B by \textbf{1.2\x} on A100, \textbf{1.4\x} on L40S; and Qwen1.5-72B by \textbf{2.4\x} on A100, \textbf{3.5\x} on L40S, compared to TensorRT-LLM. Remarkably, \system on L40S GPU can achieve even higher throughput than TensorRT-LLM on A100. Code is released at \url{https://github.com/mit-han-lab/omniserve}. %

\end{abstract}
